\subsection{Conventional Denoising Baselines}
\label{sec:conventional-baselines}

Classical denoising methods form the standard baseline for laser ultrasonic signal enhancement. Wiener filtering estimates the signal via linear minimum mean square error (MMSE), assuming stationary statistics; however, its frequency-domain smoothing destroys high-frequency transient components that carry defect-indicative information. Wavelet-based denoising (DWT) decomposes the signal into multi-resolution coefficients and applies soft or hard thresholding to suppress noise \cite{chen2019wavelet}. The \texttt{db4} wavelet is commonly adopted for laser ultrasonic signals due to its favorable time-frequency localization \cite{chen2019wavelet}. BM3D extends wavelet thresholding by grouping similar signal patches into 3-D blocks and applying collaborative filtering, achieving state-of-the-art performance on aggregate noise metrics \cite{dabov2007bm3d}.

A fundamental limitation unites these approaches: all optimize the same objective functions---MSE or SNR improvement---without explicit constraints on acoustic feature preservation. For laser ultrasonic defect detection, diagnostically relevant features include arrival times of Lamb wave modes, amplitude ratios between modes, and waveform morphology. These features are not explicitly penalized during denoising, and aggressive noise suppression routinely degrades them. Chen \textit{et al.} noted that wavelet thresholding can distort arrival times even when SNR improves \cite{chen2019wavelet}, and BM3D's patch-based processing introduces additional phase distortion.

Consequently, no conventional method explicitly preserves features for laser ultrasonic defect detection. This gap motivates physics-constrained approaches that treat acoustic feature preservation as a first-class optimization objective.